\begin{sol}[Péndulo elástico]

\subsubsection*{Pregunta 1: Grados de libertad}

A diferencia del péndulo simple, el resorte permite variar su longitud, por lo que el sistema tiene $2$ grados de libertad. Las coordenadas generalizadas naturales son $q_1 = r$ (longitud instantánea del resorte) y $q_2 = \theta$ (ángulo respecto a la vertical).

\subsubsection*{Pregunta 2: Lagrangiano}

En coordenadas polares centradas en $O$: $x = r\sin\theta$, $y=-r\cos\theta$. La rapidez al cuadrado en polares es el resultado estándar:
$$v^2 = \dot{r}^2 + r^2\dot\theta^2$$

Energía cinética:
$$T = \frac{1}{2}m\dot{r}^2 + \frac{1}{2}mr^2\dot\theta^2$$

Energía potencial: gravitatoria más elástica, con referencia en $O$:
$$V = -mgr\cos\theta + \frac{1}{2}k(r-\ell_0)^2$$

$$\boxed{L = \frac{1}{2}m\dot{r}^2 + \frac{1}{2}mr^2\dot\theta^2 + mgr\cos\theta - \frac{1}{2}k(r-\ell_0)^2}$$

\subsubsection*{Pregunta 3: Ecuaciones de movimiento}

\textbf{Respecto a $r$:}
$$\frac{\partial L}{\partial\dot{r}} = m\dot{r}, \qquad \frac{\partial L}{\partial r} = mr\dot\theta^2 + mg\cos\theta - k(r-\ell_0)$$

$$m\ddot{r} - mr\dot\theta^2 - mg\cos\theta + k(r-\ell_0) = 0$$

$$\boxed{\ddot{r} - r\dot\theta^2 - g\cos\theta + \frac{k}{m}(r-\ell_0) = 0}$$

\textbf{Respecto a $\theta$:}
$$\frac{\partial L}{\partial\dot\theta} = mr^2\dot\theta \quad\Rightarrow\quad \frac{d}{dt}\left(\frac{\partial L}{\partial\dot\theta}\right) = mr^2\ddot\theta + 2mr\dot{r}\dot\theta$$

$$\frac{\partial L}{\partial\theta} = -mgr\sin\theta$$

$$mr^2\ddot\theta + 2mr\dot{r}\dot\theta + mgr\sin\theta = 0$$

Dividiendo por $mr$ (con $r\neq0$):

$$\boxed{r\ddot\theta + 2\dot{r}\dot\theta + g\sin\theta = 0}$$

El término $2\dot r\dot\theta$ es la contribución tipo Coriolis: acopla el estiramiento radial con la oscilación angular, siendo responsable de fenómenos como la transferencia de energía entre el modo radial y el modo angular (resonancia $1{:}2$) característica del péndulo elástico.

\subsubsection*{Pregunta 4: Límite de resorte infinitamente rígido}

Si $k\to\infty$, el resorte se vuelve efectivamente inextensible y $r(t) \to \ell_0 = \text{constante}$, de modo que $\dot{r}=\ddot{r}=0$. Sustituyendo en la ecuación para $\theta$:

$$\ell_0\ddot\theta + 2(0)\dot\theta + g\sin\theta = 0$$

$$\boxed{\ddot\theta + \frac{g}{\ell_0}\sin\theta = 0}$$

que es exactamente la ecuación del péndulo simple de longitud $\ell_0$ (comparar con el Problema 1 de este tema), confirmando la consistencia del modelo: el péndulo elástico se reduce al péndulo simple cuando se elimina el grado de libertad radial.

\end{sol}
